\documentclass[10pt,a4paper]{amsart}
\usepackage{amsmath,amssymb,amsthm,graphicx}
% Retained minimal preamble for the Stacks Project chapter flat.tex (More on Flatness).
% Verbatim excerpt of the chapter preamble: Xy-pic setup (native preamble lines 14-18) and the native
% theorem declarations with the mathematical macros (native preamble lines 154-217). Copied mechanically;
% nothing added, removed or redefined. The wrapper class and the amsmath/amssymb/amsthm/graphicx base
% packages are supplied by the assembler, so the class/externaldocument part of the native preamble is omitted.
% For commutative diagrams we use Xy-pic
\usepackage[all]{xy}

% We use 2cell for 2-commutative diagrams.
\xyoption{2cell}
% Theorem environments.
%
\theoremstyle{plain}
\newtheorem{theorem}[subsection]{Theorem}
\newtheorem{proposition}[subsection]{Proposition}
\newtheorem{lemma}[subsection]{Lemma}

\theoremstyle{definition}
\newtheorem{definition}[subsection]{Definition}
\newtheorem{example}[subsection]{Example}
\newtheorem{exercise}[subsection]{Exercise}
\newtheorem{situation}[subsection]{Situation}

\theoremstyle{remark}
\newtheorem{remark}[subsection]{Remark}
\newtheorem{remarks}[subsection]{Remarks}

\numberwithin{equation}{subsection}

% Macros
%
\def\lim{\mathop{\mathrm{lim}}\nolimits}
\def\colim{\mathop{\mathrm{colim}}\nolimits}
\def\Spec{\mathop{\mathrm{Spec}}}
\def\Hom{\mathop{\mathrm{Hom}}\nolimits}
\def\Ext{\mathop{\mathrm{Ext}}\nolimits}
\def\SheafHom{\mathop{\mathcal{H}\!\mathit{om}}\nolimits}
\def\SheafExt{\mathop{\mathcal{E}\!\mathit{xt}}\nolimits}
\def\Sch{\mathit{Sch}}
\def\Mor{\mathop{\mathrm{Mor}}\nolimits}
\def\Ob{\mathop{\mathrm{Ob}}\nolimits}
\def\Sh{\mathop{\mathit{Sh}}\nolimits}
\def\NL{\mathop{N\!L}\nolimits}
\def\CH{\mathop{\mathrm{CH}}\nolimits}
\def\proetale{{pro\text{-}\acute{e}tale}}
\def\etale{{\acute{e}tale}}
\def\QCoh{\mathit{QCoh}}
\def\Ker{\mathop{\mathrm{Ker}}}
\def\Im{\mathop{\mathrm{Im}}}
\def\Coker{\mathop{\mathrm{Coker}}}
\def\Coim{\mathop{\mathrm{Coim}}}

% Boxtimes
%
\DeclareMathSymbol{\boxtimes}{\mathbin}{AMSa}{"02}

%
% Macros for moduli stacks/spaces
%
\def\QCohstack{\mathcal{QC}\!\mathit{oh}}
\def\Cohstack{\mathcal{C}\!\mathit{oh}}
\def\Spacesstack{\mathcal{S}\!\mathit{paces}}
\def\Quotfunctor{\mathrm{Quot}}
\def\Hilbfunctor{\mathrm{Hilb}}
\def\Curvesstack{\mathcal{C}\!\mathit{urves}}
\def\Polarizedstack{\mathcal{P}\!\mathit{olarized}}
\def\Complexesstack{\mathcal{C}\!\mathit{omplexes}}
% \Pic is the operator that assigns to X its picard group, usage \Pic(X)
% \Picardstack_{X/B} denotes the Picard stack of X over B
% \Picardfunctor_{X/B} denotes the Picard functor of X over B
\def\Pic{\mathop{\mathrm{Pic}}\nolimits}
\def\Picardstack{\mathcal{P}\!\mathit{ic}}
\def\Picardfunctor{\mathrm{Pic}}
\def\Deformationcategory{\mathcal{D}\!\mathit{ef}}

% --- appended 2026-10-12 (Chief): the project own definitions of the reference/slogan/history
% display environments, copied mechanically from the shared Stacks preamble lines 8-12
% (\usepackage{verbatim} plus the three \newenvironment lines, which the chapter preamble keeps
% verbatim from the project preamble). Nothing added, removed or redefined: these environments
% carry no mathematical content and the project itself typesets them as comments. Without these
% three lines a counted statement whose own text uses them fails with Environment slogan undefined.
\usepackage{verbatim}
\newenvironment{reference}{\comment}{\endcomment}
\newenvironment{slogan}{\comment}{\endcomment}
\newenvironment{history}{\comment}{\endcomment}
\usepackage{cleveref}
\crefname{originallemma}{original lemma}{original lemmas}
\Crefname{originallemma}{Original lemma}{Original lemmas}
\makeatletter
\newlabel{algebra-lemma-Noetherian-coherent}{{05CY}{source}{}{sourcepointer-1}{}}
\newlabel{algebra-lemma-Noetherian-coherent@cref}{{[originallemma][0][]05CY}{[1][0][]source}{}{}{}}
\makeatother
\crefname{originallemma}{original lemma}{original lemmas}
\Crefname{originallemma}{Original lemma}{Original lemmas}
\makeatletter
\newlabel{algebra-lemma-product-over-Noetherian-ring}{{05D0}{source}{}{sourcepointer-2}{}}
\newlabel{algebra-lemma-product-over-Noetherian-ring@cref}{{[originallemma][0][]05D0}{[1][0][]source}{}{}{}}
\makeatother
\crefname{originallemma}{original lemma}{original lemmas}
\Crefname{originallemma}{Original lemma}{Original lemmas}
\makeatletter
\newlabel{algebra-lemma-pure-submodule-ML}{{059N}{source}{}{sourcepointer-3}{}}
\newlabel{algebra-lemma-pure-submodule-ML@cref}{{[originallemma][0][]059N}{[1][0][]source}{}{}{}}
\makeatother
\crefname{originallemma}{original lemma}{original lemmas}
\Crefname{originallemma}{Original lemma}{Original lemmas}
\makeatletter
\newlabel{algebra-lemma-ui-flat-domain}{{058P}{source}{}{sourcepointer-4}{}}
\newlabel{algebra-lemma-ui-flat-domain@cref}{{[originallemma][0][]058P}{[1][0][]source}{}{}{}}
\makeatother
\crefname{originalproposition}{original proposition}{original propositions}
\Crefname{originalproposition}{Original proposition}{Original propositions}
\makeatletter
\newlabel{algebra-proposition-characterize-coherent}{{05CZ}{source}{}{sourcepointer-5}{}}
\newlabel{algebra-proposition-characterize-coherent@cref}{{[originalproposition][0][]05CZ}{[1][0][]source}{}{}{}}
\makeatother
\crefname{originallemma}{original lemma}{original lemmas}
\Crefname{originallemma}{Original lemma}{Original lemmas}
\makeatletter
\newlabel{more-algebra-lemma-ui-completion-direct-sum-into-product}{{05BC}{source}{}{sourcepointer-6}{}}
\newlabel{more-algebra-lemma-ui-completion-direct-sum-into-product@cref}{{[originallemma][0][]05BC}{[1][0][]source}{}{}{}}
\makeatother
\title{The completion of a direct sum of copies of a complete Noetherian ring is flat and Mittag-Leffler}
\author{The Stacks Project Authors}
\date{stacks-project:flat.tex:local-snapshot-20261011; GNU FDL 1.2 or later}
\begin{document}
\maketitle
\noindent\textit{Editorial note.} All delivered mathematical text is the author's own, copied from the pinned Stacks Project chapter file flat.tex; no statement or proof step was rewritten, reordered or removed.\par\medskip
\noindent\textit{Editorial note.} Preamble provenance: the chapter's own retained minimal preamble (sources/stacks-flat/retained-preamble.tex, lines 1-85) is a verbatim excerpt of the native Xy-pic setup and of the native theorem declarations and mathematical macros, followed by the project's own reference/slogan/history display-environment definitions copied mechanically from the shared preamble; the wrapper class and the amsmath/amssymb/amsthm/graphicx base packages are supplied by the assembler, and extra\_packages is left empty because everything needed is already in the retained preamble. No macro was invented or redefined.\par\medskip
\noindent\textit{Editorial note.} Same-chapter pointers are located at their \textbackslash{}label lines inside the pinned chapter file; cross-chapter pointers are located in the project's own tag list staged in the snapshot (sources/stacks-chow/tags.txt), at the verbatim lines carrying their canonical tags.\par\medskip
\noindent\textit{Allowed standard prerequisites (editorial).}\begin{itemize}
\item Universal injectivity of the completion map into the product, flatness of the cokernel of a universally injective map, and the Mittag-Leffler criterion for a pure submodule. More on Algebra, label more-algebra-lemma-ui-completion-direct-sum-into-product (proof line 2013); Algebra, labels algebra-lemma-ui-flat-domain and algebra-lemma-pure-submodule-ML (proof lines 2017-2018). Shows that the completion is a pure submodule of the product of copies of the ring, so that flatness and the Mittag-Leffler property descend from the product.
\item Coherence and flatness of a product of copies of a Noetherian ring, and its Mittag-Leffler property. Algebra, labels algebra-proposition-characterize-coherent (proof line 2021), algebra-lemma-Noetherian-coherent (proof line 2023) and algebra-lemma-product-over-Noetherian-ring (proof line 2025). Gives that the product of copies of a Noetherian ring is flat and Mittag-Leffler, which combined with the purity statement gives the conclusion.
\end{itemize}\medskip
\ifdefined\hypertarget\hypertarget{sourcepointer-1}{}\fi
\noindent\textit{Source reference.} Original lemma 05CY: sources/stacks-chow/tags.txt, source lines 6936--6936. Navigation pointer for \textbackslash{}ref\{algebra-lemma-Noetherian-coherent\} used inside the counted span (line 2023); the argument that consumes it is retained in the span.\par\medskip
\ifdefined\hypertarget\hypertarget{sourcepointer-2}{}\fi
\noindent\textit{Source reference.} Original lemma 05D0: sources/stacks-chow/tags.txt, source lines 6938--6938. Navigation pointer for \textbackslash{}ref\{algebra-lemma-product-over-Noetherian-ring\} used inside the counted span (line 2026); the argument that consumes it is retained in the span.\par\medskip
\ifdefined\hypertarget\hypertarget{sourcepointer-3}{}\fi
\noindent\textit{Source reference.} Original lemma 059N: sources/stacks-chow/tags.txt, source lines 6821--6821. Navigation pointer for \textbackslash{}ref\{algebra-lemma-pure-submodule-ML\} used inside the counted span (line 2018); the argument that consumes it is retained in the span.\par\medskip
\ifdefined\hypertarget\hypertarget{sourcepointer-4}{}\fi
\noindent\textit{Source reference.} Original lemma 058P: sources/stacks-chow/tags.txt, source lines 6787--6787. Navigation pointer for \textbackslash{}ref\{algebra-lemma-ui-flat-domain\} used inside the counted span (line 2017); the argument that consumes it is retained in the span.\par\medskip
\ifdefined\hypertarget\hypertarget{sourcepointer-5}{}\fi
\noindent\textit{Source reference.} Original proposition 05CZ: sources/stacks-chow/tags.txt, source lines 6937--6937. Navigation pointer for \textbackslash{}ref\{algebra-proposition-characterize-coherent\} used inside the counted span (line 2021); the argument that consumes it is retained in the span.\par\medskip
\ifdefined\hypertarget\hypertarget{sourcepointer-6}{}\fi
\noindent\textit{Source reference.} Original lemma 05BC: sources/stacks-chow/tags.txt, source lines 6880--6880. Navigation pointer for \textbackslash{}ref\{more-algebra-lemma-ui-completion-direct-sum-into-product\} used inside the counted span (line 2013); the argument that consumes it is retained in the span.\par\medskip
\begin{lemma}
\label{lemma-completed-direct-sum-ML}
Let $R$ be a ring. Let $I \subset R$ be an ideal. Let $A$ be a set.
Assume $R$ is Noetherian and complete with respect to $I$. The completion
$(\bigoplus\nolimits_{\alpha \in A} R)^\wedge$
is flat and Mittag-Leffler.
\end{lemma}

\begin{proof}
By
More on Algebra, Lemma
\ref{more-algebra-lemma-ui-completion-direct-sum-into-product}
the map $(\bigoplus\nolimits_{\alpha \in A} R)^\wedge
\to \prod_{\alpha \in A} R$ is universally injective.
Thus, by
Algebra, Lemmas \ref{algebra-lemma-ui-flat-domain} and
\ref{algebra-lemma-pure-submodule-ML}
it suffices to show that $\prod_{\alpha \in A} R$ is flat and Mittag-Leffler.
By
Algebra, Proposition \ref{algebra-proposition-characterize-coherent}
(and
Algebra, Lemma \ref{algebra-lemma-Noetherian-coherent})
we see that $\prod_{\alpha \in A} R$ is flat.
Thus we conclude because a product of copies of $R$ is Mittag-Leffler, see
Algebra, Lemma \ref{algebra-lemma-product-over-Noetherian-ring}.
\end{proof}

\end{document}
